% Proof proposal; full raw zeros and arbitrary fixed coefficients retained.
\section{Proof of the ordinary-conductor refinement}\label{mcs:details}
\subsection{Primitive Gauss sums at every conductor}
Let $\mathcal O=\mathbb Z[\omega]$, $\delta=1+2\omega$, and choose one
generator $f$ of each nonzero ideal $\mathfrak f$. Write
$e_f(x)=\exp(2\pi i\operatorname{Tr}(x/(\delta f)))$,
$U(\mathfrak f)=(\mathcal O/\mathfrak f)^\times$, and
$\varphi(\mathfrak f)=|U(\mathfrak f)|$.
This additive character is primitive, since the different is $(\delta)$.
For a primitive multiplicative character $\rho$ modulo $\mathfrak f$,
extended by zero at nonunits, set
$\tau_f(\rho)=\sum_{x\bmod\mathfrak f}\rho(x)e_f(x)$.
We verify the full-modulus identity
\begin{equation}\label{mcs:gauss}
 |\tau_f(\rho)|^2=N\mathfrak f,\qquad
 \sum_{x\bmod\mathfrak f}\rho(x)e_f(zx)
       =\overline{\rho(z)}\tau_f(\rho)\quad(z\in\mathcal O).
\end{equation}
No squarefree-conductor hypothesis is used, so this includes wild
prime powers in the fixed $S$-part.

The Ramanujan sum is
\[
 c_{\mathfrak f}(y)=\sum_{x\in U(\mathfrak f)}e_f(xy)
   =\sum_{\mathfrak d\mid(\mathfrak f,(y))}
                   N\mathfrak d\,\mu(\mathfrak f/\mathfrak d).
\]
Indeed insert $\mathbf1_{(x,\mathfrak f)=1}
=\sum_{\mathfrak b\mid(x,\mathfrak f)}\mu(\mathfrak b)$;
the complete additive sum with $\mathfrak b\mid x$ is
$N(\mathfrak f/\mathfrak b)$ when
$\mathfrak f/\mathfrak b\mid(y)$ and zero otherwise.
Changing variables $x=ay$ in the Gauss square now gives
\begin{align*}
 |\tau_f(\rho)|^2
 &=\sum_{a\in U(\mathfrak f)}\rho(a)c_{\mathfrak f}(a-1)\\
 &=\sum_{\mathfrak d\mid\mathfrak f}N\mathfrak d\,
       \mu(\mathfrak f/\mathfrak d)
       \sum_{\substack{a\in U(\mathfrak f)\\a\equiv1\pmod{\mathfrak d}}}
                  \rho(a)=N\mathfrak f.
\end{align*}
For every proper $\mathfrak d\mid\mathfrak f$, the reduction
$U(\mathfrak f)\to U(\mathfrak d)$ is surjective by local lifting and CRT.
Primitivity says $\rho$ is nontrivial on its kernel, so that inner sum
vanishes. For $\mathfrak d=\mathfrak f$ it equals one.
For unit $z$, substitution proves the second identity in~\eqref{mcs:gauss}.
For all frequencies together, finite additive Parseval gives
\[
 \sum_{z\bmod\mathfrak f}
   \left|\sum_x\rho(x)e_f(zx)\right|^2
       =N\mathfrak f\,\varphi(\mathfrak f).
\]
The unit frequencies already exhaust this quantity by the first identity;
every nonunit-frequency sum is therefore zero. This proves the asserted
zero extension. For $\mathfrak f=\mathcal O$ there is one residue class
and one character; take $\tau_1=1$ and the single frequency zero.

\subsection{The classical ordinary sieve in these lattice conventions}
For $D,Q\ge1$, coefficients $z_\alpha$ supported on any subset of
$\{\alpha\in\mathcal O:|\alpha|^2\le C_0D\}$ satisfy
\begin{equation}\label{mcs:ordinary}
 \sum_{N\mathfrak f\le Q}\ \sum_{\rho\ \mathrm{primitive}\bmod\mathfrak f}
 \left|\sum_\alpha z_\alpha\rho(\alpha)\right|^2
 \ll_{\mathcal O,C_0}(D+Q^2)\sum_\alpha|z_\alpha|^2.
\end{equation}
Here the primitive characters may be restricted to the unit-trivial
subfamily. We give the standard proof to make the conductor comparison
explicit.

First the reduced fractions $x/f$, with $x\in U(\mathfrak f)$ and
$N\mathfrak f\le Q$, are distinct modulo $\mathcal O$.
Their denominator ideal is intrinsic:
$\{a\in\mathcal O:a(x/f)\in\mathcal O\}=\mathfrak f$.
Thus equality modulo $\mathcal O$ forces the same denominator ideal
and residue numerator. For different fractions and every $k\in\mathcal O$,
\[
 \left|\frac{x}{f}-\frac{y}{g}-k\right|
   =\frac{|xg-yf-kfg|}{|fg|}\ge\frac1Q,
\]
because the numerator is a nonzero Eisenstein integer and
$|fg|=(N\mathfrak f N\mathfrak g)^{1/2}\le Q$.
The fixed codifferent factor $1/\delta$ changes only the lattice constants.
The modulus-one frequency zero obeys the same separation from all others.

For completeness, a $\rho_0$-separated set $t_r$ in a fixed
two-dimensional torus obeys the additive bound
\begin{equation}\label{mcs:additive}
 \sum_r\left|\sum_{|\alpha|\ll\sqrt D}
        z_\alpha\exp(2\pi i\operatorname{Tr}(\alpha t_r))\right|^2
       \ll(D+\rho_0^{-2})\sum_\alpha|z_\alpha|^2.
\end{equation}
Use fixed linear coordinates to write the lattice as $\mathbb Z^2$ and
its additive pairing as $e(n\cdot t)$. The support lies in
$|n_i|\le T$ with $T\asymp\max(1,\sqrt D)$.
To prove the dual inequality, majorize this box indicator by
$e^{2\pi}\exp(-\pi\|n\|^2/T^2)$, expand the square, and apply Gaussian
Poisson summation. Up to fixed constants the resulting kernel is
\[
 T^2\sum_{\ell\in\mathbb Z^2}
       \exp(-cT^2\|t_r-t_s+\ell\|^2).
\]
Take $\rho_0$ at most a fixed torus constant. The lifted points are
$c'\rho_0$-separated, so disjoint disks bound their count in a disk
of radius $L$ by $O(1+(L/\rho_0)^2)$.
Gaussian annular summation bounds every absolute kernel row sum by
$O(T^2+\rho_0^{-2})$. Schur's inequality and finite-dimensional duality
prove~\eqref{mcs:additive}, with constants independent of the number
of frequencies.

Put $\mathcal S(x/f)=\sum_\alpha z_\alpha e_f(\alpha x)$.
By~\eqref{mcs:gauss},
\[
 \rho(\alpha)=\tau_f(\overline\rho)^{-1}
       \sum_{x\in U(\mathfrak f)}\overline{\rho(x)}e_f(\alpha x).
\]
After taking squared absolute values, enlarging the primitive-character
sum to all characters of $U(\mathfrak f)$ and using orthogonality yields
\[
 \sum_{\rho\ \mathrm{primitive}\bmod\mathfrak f}
       \left|\sum_\alpha z_\alpha\rho(\alpha)\right|^2
 \le\frac{\varphi(\mathfrak f)}{N\mathfrak f}
       \sum_{x\in U(\mathfrak f)}|\mathcal S(x/f)|^2
 \le\sum_{x\in U(\mathfrak f)}|\mathcal S(x/f)|^2.
\]
Apply~\eqref{mcs:additive} with $\rho_0\gg Q^{-1}$ to prove
\eqref{mcs:ordinary}.

Finally $\mathcal O$ has class number one. Choose one generator
$\alpha_m$ of each ideal column, with $|\alpha_m|^2=Nm$; these are an
injective subset of the lattice disk. A primitive trivial-infinity-type
Hecke character is a unit-trivial primitive residue character on
$\mathcal O/\mathfrak f$, and its ideal value equals its value on
$\alpha_m$. Descent of the residue character to a proper divisor would
preserve unit-triviality and contradict primitive Hecke conductor.
Thus~\eqref{mcs:ordinary} applies to arbitrary ideal coefficients of
norm at most $2D$, including squarefree ideals of nonsquarefree rational norm.

\subsection{Raw zeros, tuple multiplicity and block transfer}
Freeze $r_6$ and the literal unit/$S$-factor, putting
\[
 c_m=b_m\chi_m(u_0)\mathbf1_{(m,r_6)=1},\qquad
       \sum_m|c_m|^2\le\sum_m|b_m|^2.
\]
This is independent of every first-five variable. In particular the
arbitrary overlapping sixth-power factor is not replaced by one.
Let $v=\prod_{j=1}^5r_j^j$ denote the chosen literal product.
Every outside-$S$ valuation is between one and five.
BGL, \href{https://arxiv.org/pdf/1112.1650v1}{arXiv:1112.1650v1},
Lemma~2.1, together with the established finite $S$-unit comparison
in~\eqref{wc:literal-reciprocity}, gives its primitive literal Kummer
character $\kappa_v$ of trivial infinity type, with
\[
 q=r_1\cdots r_5\mid\operatorname{cond}(\kappa_v)\mid\mathfrak c_Sq,
 \qquad \kappa_v(m)=\chi_m(v)\quad((m,S)=1).
\]
The equality includes collisions with $q$: both sides are zero.
The fixed $\mathfrak c_S$ includes all wild conductor types and literal
comparison twists. No row-dependent phase enters the coefficients.

A fixed primitive character determines its outside-$S$ conductor $q$.
Each prime of $q$ has at most five assignments to the pairwise-coprime
$r_1,\ldots,r_5$, so tuple multiplicity is at most
$5^{\omega(q)}\ll_\eta(Nq)^\eta$.
All conductors in the dyadic block have norm at most $C_SQ$.
Enlarge the positive row sum, apply~\eqref{mcs:ordinary}, and count
$O(A_6)$ possibilities for the frozen arbitrary $r_6$.
This proves~\eqref{mcs:block-bound}; the actual capacity, squarefree,
ray and inter-factor restrictions only shrink the positive sum.
Equivalently $\mathcal E_A(b)\le\Lambda_A\sum_m|b_m|^2$;
there is no coefficient norm factor in the definition of $\Lambda_A$.

\subsection{Uniform coefficient and principal-residue assembly}
The other block estimate is exactly the accepted powered-BGL input
\eqref{wc:block}, whose squarefree powers are justified in BGL
Section~3.2. The capacity comparisons and elementary minimum in
Section~\ref{mcs:core} give
\[
 \Lambda_A\ll(HD)^\varepsilon
          \{H+DH^{1/6}+(DH)^{2/3}\}.
\]
Sum logarithmically many blocks with smaller internal exponents.
An $S$-supported literal factor $\mathfrak b$ reduces capacity to
$H/N\mathfrak b$; the required sums over $\mathfrak b$ converge
at powers $1$, $1/6$, $2/3$. Literal units have multiplicity at most six.
This proves the raw bound with arbitrary fixed $b_m$.

For the actual coefficients use~\eqref{wc:l2}:
$|A_U(m)|\le3^{\omega(m)}\ll_\eta(Nm)^\eta$ gives
\[
 b_m=D^{-1/2}A_U(m)\nu(m)W(Nm/D),\qquad
       \sum_m|b_m|^2\ll_\eta\|W\|_\infty^2D^\eta,
\]
uniformly for $U\ge1$, $U^2\le D$. The common $W$ is not optimized
separately for different rows. If the primitive $\psi_u$ is principal,
$\nu$ is unramified outside $S$, so tame Kummer inertia forces every
outside-$S$ valuation of $u$ to be divisible by six. Ideal counting and
the convergent $S$-supported sum at power $1/6$ give
$O_{\nu,S}(H^{1/6})$ such literal rows.
In the exact formula~\eqref{wc:principal}, $|\rho_T|\ll_K1$,
$|\widehat W(1)|\le\|W\|_\infty$, and
\[
 \sum_d\frac{|\beta_U(d)|}{Nd}
 \le\left(\sum_{Na\le U}\frac1{Na}\right)^2
 \ll\log^2(2U).
\]
Thus $|P_u|\ll\sqrt D\|W\|_\infty\log^2(2U)$ on principal rows,
and
\[
 \sum_{0<Nu\le H}|P_u|^2
       \ll DH^{1/6}\|W\|_\infty^2\log^4(2U).
\]
This fits the new envelope. The exact identity $E_u=I_u-P_u$ and
$|E_u|^2\le2|I_u|^2+2|P_u|^2$, with smaller internal exponents,
prove~\eqref{mcs:uniform}. No cutoff limit or profile supremum occurs.

The ordinary comparison is an application refinement. A direct BGL
radical relabelling would additionally require reciprocity of its mixed
family, not just its conductor count: for pairwise-coprime $a,b,a',b'$,
the quotient after the original fixed ray phases is
$\vartheta_b(a')/\vartheta_a(b')$, which varies outside $S$.
GL, \href{https://arxiv.org/pdf/1112.1642v2}{arXiv:1112.1642v2},
Definition~1(2), is therefore not supplied by that shortcut.
This does not rule out a different mixed-family theorem; it is simply
unneeded for the proved weaker ordinary comparison.
The remaining exponent gap and every independent signed, spectral,
formal and GRH obligation stated in Section~\ref{mcs:core} remain open.
